% GENERATED FILE. Edit canonical lessons, then run npm run generate:books.
% canonical-source-hash: fnv1a64:a869a40f2e4d1f3a
% canonical-lessons: PT-C407-capo, PT-C407-bagageira, PT-C407-pisca, PT-C407-parafuso, PT-C407-desaparafusar

\chapter{Debaixo do capô}
\label{ch:pt-debaixo-do-capo}
\hlchaptermodality{407}

\begin{chapteropening}
\textbf{By the end of this chapter:} \emph{I can talk about parts of a car: the bonnet, the boot, the indicator, a screw and unscrewing it.}

Complete the last lesson of chapter 407: I can talk about parts of a car: the bonnet, the boot, the indicator, a screw and unscrewing it.
\end{chapteropening}

\section[o capô]{o capô — the bonnet (of a car)}

\label{lesson:PT-C407-capo}



\begin{quote}
Before the new one: say the Portuguese for the tow truck, then the Portuguese for the repair garage.
\end{quote}

\subsection*{You'll want to know: o capô}
\textbf{o capô} — \textquotedblleft{}the bonnet (of a car)\textquotedblright{}.

\begin{quote}
O mecânico abriu o capô para ver o motor.
\end{quote}

\begin{quote}
\emph{The mechanic opened the bonnet to check the engine.}
\end{quote}

In Brazil also \emph{o capô}; in the US, the hood.

\begin{cousinweb}
\emph{capô} was borrowed from French \emph{capot} 'bonnet, cover' (EtymDB).
\end{cousinweb}

\begin{grammarlens}[title={What you've built}]
One more word for telling, explaining and describing at B1.
\end{grammarlens}

\subsection*{Your turn}
\begin{itemize}
\raggedright
  \item \emph{Say it:} \emph{o capô}
  \item \emph{Say it:} \emph{o capô}, once more
  \item \emph{Say it:} a sentence of your own with \emph{o capô}
  \item \emph{Recall:} say the Portuguese for the tow truck, then the Portuguese for the repair garage, then say \emph{o capô} again
\end{itemize}

\begin{tcolorbox}[breakable,colback=teal!4,colframe=teal!35!black,title={Before you move on}]
What does o capô mean? (\textquotedblleft{}The bonnet (of a car)\textquotedblright{}.) Say it once more.
\end{tcolorbox}
\section[a bagageira]{a bagageira — the boot (of a car)}

\label{lesson:PT-C407-bagageira}



\begin{quote}
Before the new one: say the Portuguese for the repair garage, then the Portuguese for the bonnet.
\end{quote}

\subsection*{You'll want to know: a bagageira}
\textbf{a bagageira} — \textquotedblleft{}the boot (of a car)\textquotedblright{}.

\begin{quote}
As malas grandes não cabem na bagageira.
\end{quote}

\begin{quote}
\emph{The big suitcases do not fit in the boot.}
\end{quote}

In Brazil: \emph{o porta-malas}. In Portugal people also say \emph{a mala do carro}.

\begin{cousinweb}
\emph{bagageira} is built on \emph{bagagem} 'luggage' with \emph{-eira}, a place for things; \emph{bagagem} came from French \emph{bagage} (EtymDB lists the French word).
\end{cousinweb}

\begin{grammarlens}[title={What you've built}]
One more word for telling, explaining and describing at B1.
\end{grammarlens}

\subsection*{Your turn}
\begin{itemize}
\raggedright
  \item \emph{Say it:} \emph{a bagageira}
  \item \emph{Say it:} \emph{a bagageira}, once more
  \item \emph{Say it:} a sentence of your own with \emph{a bagageira}
  \item \emph{Recall:} say the Portuguese for the repair garage, then the Portuguese for the bonnet, then say \emph{a bagageira} again
\end{itemize}

\begin{tcolorbox}[breakable,colback=teal!4,colframe=teal!35!black,title={Before you move on}]
What does a bagageira mean? (\textquotedblleft{}The boot (of a car)\textquotedblright{}.) Say it once more.
\end{tcolorbox}
\section[o pisca]{o pisca — the indicator, the turn signal}

\label{lesson:PT-C407-pisca}



\begin{quote}
Before the new one: say the Portuguese for the bonnet, then the Portuguese for the boot.
\end{quote}

\subsection*{You'll want to know: o pisca}
\textbf{o pisca} — \textquotedblleft{}the indicator, the turn signal\textquotedblright{}.

\begin{quote}
Ele nunca liga o pisca antes de virar.
\end{quote}

\begin{quote}
\emph{He never puts the indicator on before turning.}
\end{quote}

Short for \emph{o pisca-pisca}. \emph{Piscar o olho} — to wink.

\begin{cousinweb}
\emph{pisca} is built on \emph{piscar} 'to blink'. The origin of \emph{piscar} is not settled; EtymDB has no entry for it.
\end{cousinweb}

\begin{grammarlens}[title={What you've built}]
One more word for telling, explaining and describing at B1.
\end{grammarlens}

\subsection*{Your turn}
\begin{itemize}
\raggedright
  \item \emph{Say it:} \emph{o pisca}
  \item \emph{Say it:} \emph{o pisca}, once more
  \item \emph{Say it:} a sentence of your own with \emph{o pisca}
  \item \emph{Recall:} say the Portuguese for the bonnet, then the Portuguese for the boot, then say \emph{o pisca} again
\end{itemize}

\begin{tcolorbox}[breakable,colback=teal!4,colframe=teal!35!black,title={Before you move on}]
What does o pisca mean? (\textquotedblleft{}The indicator, the turn signal\textquotedblright{}.) Say it once more.
\end{tcolorbox}
\section[o parafuso]{o parafuso — the screw}

\label{lesson:PT-C407-parafuso}



\begin{quote}
Before the new one: say the Portuguese for the boot, then the Portuguese for the indicator.
\end{quote}

\subsection*{You'll want to know: o parafuso}
\textbf{o parafuso} — \textquotedblleft{}the screw\textquotedblright{}.

\begin{quote}
Falta um parafuso na porta do carro.
\end{quote}

\begin{quote}
\emph{There is a screw missing from the car door.}
\end{quote}

\emph{Ter um parafuso a menos} — to have a screw loose.

\begin{cousinweb}
The origin of \emph{parafuso} is not settled; it is sometimes linked with \emph{fuso} 'spindle', from Latin \emph{fūsus} (Lewis \& Short), but EtymDB has no entry for it.
\end{cousinweb}

\begin{grammarlens}[title={What you've built}]
One more word for telling, explaining and describing at B1.
\end{grammarlens}

\subsection*{Your turn}
\begin{itemize}
\raggedright
  \item \emph{Say it:} \emph{o parafuso}
  \item \emph{Say it:} \emph{o parafuso}, once more
  \item \emph{Say it:} a sentence of your own with \emph{o parafuso}
  \item \emph{Recall:} say the Portuguese for the boot, then the Portuguese for the indicator, then say \emph{o parafuso} again
\end{itemize}

\begin{tcolorbox}[breakable,colback=teal!4,colframe=teal!35!black,title={Before you move on}]
What does o parafuso mean? (\textquotedblleft{}The screw\textquotedblright{}.) Say it once more.
\end{tcolorbox}
\section[desaparafusar]{desaparafusar — to unscrew}

\label{lesson:PT-C407-desaparafusar}



\begin{quote}
Before the new one: say the Portuguese for the indicator, then the Portuguese for the screw.
\end{quote}

\subsection*{You'll want to know: desaparafusar}
\textbf{desaparafusar} — \textquotedblleft{}to unscrew\textquotedblright{}.

\begin{quote}
O mecânico desaparafusou a roda com cuidado.
\end{quote}

\begin{quote}
\emph{The mechanic carefully unscrewed the wheel.}
\end{quote}

The opposite is \emph{aparafusar}, to screw in.

\begin{cousinweb}
\emph{desaparafusar} is built from \emph{des-} 'un-' and \emph{aparafusar} 'to screw in', on \emph{parafuso} 'screw', whose own origin is not settled (EtymDB lists \emph{aparafusar} without a source).
\end{cousinweb}

\begin{grammarlens}[title={What you've built}]
One more word for telling, explaining and describing at B1.
\end{grammarlens}

\subsection*{Your turn}
\begin{itemize}
\raggedright
  \item \emph{Say it:} \emph{desaparafusar}
  \item \emph{Say it:} \emph{desaparafusar}, once more
  \item \emph{Say it:} a sentence of your own with \emph{desaparafusar}
  \item \emph{Recall:} say the Portuguese for the indicator, then the Portuguese for the screw, then say \emph{desaparafusar} again
\end{itemize}

\begin{tcolorbox}[breakable,colback=teal!4,colframe=teal!35!black,title={Before you move on}]
What does desaparafusar mean? (\textquotedblleft{}To unscrew\textquotedblright{}.) Say it once more.
\end{tcolorbox}
